%%%%%%%%%%%%%%%%%%%%%%%%%%%%%%%%%%%%%%%%%
% Lachaise Assignment
% LaTeX Template
% Version 1.0 (26/6/2018)
%
% This template originates from:
% http://www.LaTeXTemplates.com
%
% Authors:
% Marion Lachaise & François Févotte
% Vel (vel@LaTeXTemplates.com)
%
% License:
% CC BY-NC-SA 3.0 (http://creativecommons.org/licenses/by-nc-sa/3.0/)
% 
%%%%%%%%%%%%%%%%%%%%%%%%%%%%%%%%%%%%%%%%%

%----------------------------------------------------------------------------------------
%	PACKAGES AND OTHER DOCUMENT CONFIGURATIONS
%----------------------------------------------------------------------------------------

\documentclass{article}
\usepackage{graphicx}
\usepackage{amsmath}
\usepackage{bm}
\usepackage{steinmetz}
\graphicspath{{images/}}
\usepackage{float} 
\usepackage{mathtools}
\usepackage[american,RPvoltages]{circuiTikz}
\usepackage{tikz, tikz-cd}
\usepackage{pgfplots}
\usepackage{siunitx}
\usetikzlibrary{spy,shapes,arrows, arrows.meta, bending}
\pgfplotsset{width=10cm,compat=1.18}
\input{structure.tex} % Include the file specifying the document structure and custom commands
%\renewcommand{\figurename}{Figura}
\usepackage{tcolorbox}
\usepackage{lipsum}
\tcbuselibrary{listings,theorems,skins,vignette,raster,breakable,magazine,poster,fitting,hooks}
\usepackage{svg}
%----------------------------------------------------------------------------------------
%	ASSIGNMENT INFORMATION
%----------------------------------------------------------------------------------------

\title{Método de superposición} % Title of the assignment
\author{Belmont Dubois} % Author name and email address

%----------------------------------------------------------------------------------------


\makeatletter
\pgfcircdeclarebipole{}{\ctikzvalof{bipoles/vsourcesin/height}}{sVpm}{\ctikzvalof{bipoles/vsourcesin/height}}{\ctikzvalof{bipoles/vsourcesin/width}}{
	
	\pgfsetlinewidth{\pgfkeysvalueof{/tikz/circuitikz/bipoles/thickness}\pgfstartlinewidth}
	\pgfpathellipse{\pgfpointorigin}{\pgfpoint{0}{\pgf@circ@res@up}}{\pgfpoint{\pgf@circ@res@left}{0}}
	\pgfusepath{draw}     
	\pgftext[bottom,rotate=90,y=\ctikzvalof{bipoles/vsourceam/margin}\pgf@circ@res@down]{\scriptsize$-$}
	\pgftext[top,rotate=90,y=\ctikzvalof{bipoles/vsourceam/margin}\pgf@circ@res@up]{\scriptsize$+$}
	
	\pgf@circ@res@up = .35\pgf@circ@res@up
	\pgfscope
	\pgftransformrotate{90}
	\pgfpathmoveto{\pgfpoint{-\pgf@circ@res@up}{0cm}}
	\pgfpathsine{\pgfpoint{.5\pgf@circ@res@up}{.4\pgf@circ@res@up}}
	\pgfpathcosine{\pgfpoint{.5\pgf@circ@res@up}{-.4\pgf@circ@res@up}}
	\pgfpathsine{\pgfpoint{.5\pgf@circ@res@up}{-.4\pgf@circ@res@up}}
	\pgfpathcosine{\pgfpoint{.5\pgf@circ@res@up}{.4\pgf@circ@res@up}}
	\pgfusepath{draw}
	\endpgfscope
}
\def\pgf@circ@sVpm@path#1{\pgf@circ@bipole@path{sVpm}{#1}}
\compattikzset{sinusoidal voltage source pm/.style = {\circuitikzbasekey, /tikz/to path=\pgf@circ@sVpm@path, \circuitikzbasekey/bipole/is voltage=true, \circuitikzbasekey/bipole/is voltageoutsideofsymbol=true, v=#1 }}
\compattikzset{sVpm/.style = {\comnpatname sinusoidal voltage source pm = #1}}

\pgfcircdeclarebipole{}{\ctikzvalof{bipoles/vsourcesin/height}}{sVpmh}{\ctikzvalof{bipoles/vsourcesin/height}}{\ctikzvalof{bipoles/vsourcesin/width}}{
	
	\pgfsetlinewidth{\pgfkeysvalueof{/tikz/circuitikz/bipoles/thickness}\pgfstartlinewidth}
	\pgfpathellipse{\pgfpointorigin}{\pgfpoint{0}{\pgf@circ@res@up}}{\pgfpoint{\pgf@circ@res@left}{0}}
	\pgfusepath{draw}     
	\pgftext[center,x=0.2\pgf@circ@res@down-\ctikzvalof{bipoles/vsourceam/margin}\pgf@circ@res@down]{\scriptsize$-$}
	\pgftext[top,rotate=90,y=\ctikzvalof{bipoles/vsourceam/margin}\pgf@circ@res@up]{\scriptsize$+$}
	
	\pgf@circ@res@up = .35\pgf@circ@res@up
	\pgfscope
	\pgftransformrotate{90}
	\pgfpathmoveto{\pgfpoint{-\pgf@circ@res@up}{0cm}}
	\pgfpathsine{\pgfpoint{.5\pgf@circ@res@up}{.4\pgf@circ@res@up}}
	\pgfpathcosine{\pgfpoint{.5\pgf@circ@res@up}{-.4\pgf@circ@res@up}}
	\pgfpathsine{\pgfpoint{.5\pgf@circ@res@up}{-.4\pgf@circ@res@up}}
	\pgfpathcosine{\pgfpoint{.5\pgf@circ@res@up}{.4\pgf@circ@res@up}}
	\pgfusepath{draw}
	\endpgfscope
}
\def\pgf@circ@sVpmh@path#1{\pgf@circ@bipole@path{sVpmh}{#1}}
\compattikzset{sinusoidal voltage source pmh/.style = {\circuitikzbasekey, /tikz/to path=\pgf@circ@sVpmh@path, \circuitikzbasekey/bipole/is voltage=true, \circuitikzbasekey/bipole/is voltageoutsideofsymbol=true, v=#1 }}
\compattikzset{sVpmh/.style = {\comnpatname sinusoidal voltage source pmh = #1}}
\makeatother

\begin{document}
	
	\maketitle % Print the title
	
	%----------------------------------------------------------------------------------------
	%	INTRODUCTION
	%----------------------------------------------------------------------------------------
	
	\section*{Análisis de circuito con superposición} % Unnumbered section
	
\begin{figure}[H]
	\centering	
	\begin{tikzpicture}[scale=0.8]
		\draw (0,0) to[R, l2=$R_{1}$ and \SI{20}{\ohm}] ++(0,4)
		to[short] ++(4,0) coordinate(n1)
		to[R, l2=$R_{2}$ and \SI{10}{\ohm}] ++(0,-4)
		to[short] ++(-4,0);
		\draw (n1) to[short, *-] ++(4,0) coordinate(n2)
		to[R, l2=$R_{3}$ and \SI{4}{\ohm}, *-*] ++(0,-4)
		to[short, -*] ++(-4,0);
		\draw (n1) to[short] ++(0,4)
		to[I, invert, l=$I_{2}$, a=\SI{3}{\ampere}] ++(4,0)
		to[short] ++(0,-4);
		\draw (n2) to[short] ++(4,0)
		to[I, invert, l2=$I_{1}$ and \SI{1}{\ampere}] ++(0,-4)
		to[short] ++(-4,0);
		\draw[->]   (2,2.5) arc(110:-110:5mm) node[midway, left, font=\footnotesize] {$I_1$};
		\draw[->]   (6,2.5) arc(110:-110:5mm) node[midway, left, font=\footnotesize] {$I_2$};
		\draw[->]   (10,2.5) arc(110:-110:5mm) node[midway, left, font=\footnotesize] {$I_3$};
		\draw[->]   (6,6.5)arc(110:-110:5mm) node[midway, left, font=\footnotesize] {$I_3$};
	\end{tikzpicture}
	\caption{}	
	\label{fig:figura1}
\end{figure}

\begin{figure}[H]
	\centering	
	\begin{tikzpicture}[scale=0.8]
		\draw (0,0) to[R, l2=$R_{1}$ and \SI{20}{\ohm}, v=$V_2$] ++(0,4)
		to[short, i=$I_1$] ++(4,0) coordinate(n1)
		to[R, l2=$R_{2}$ and \SI{10}{\ohm}, v=$V_2$, i>^=$I_a$] ++(0,-4)
		to[short, ] ++(-4,0);
		\draw (n1) to[short, i=$I_z$, *-] ++(4,0) coordinate(n2)
		to[R, l2=$R_{3}$ and \SI{4}{\ohm}, i>^=$I_b$, v=$V_3$, *-*] ++(0,-4)
		to[short, -*] ++(-4,0);
		\draw (n1) to[short, i<=$I_x$] ++(0,4)
		to[I, invert, l=$I_{X}$, a=\SI{3}{\ampere}] ++(4,0)
		to[short] ++(0,-4);
		\draw (n2) to[short] ++(4,0)
		to[I, invert, l2=$I_{1}$ and \SI{1}{\ampere}] ++(0,-4)
		to[short] ++(-4,0);
		\draw[->]   (2,2.5) arc(110:-110:5mm) node[midway, left, font=\footnotesize] {$I_1$};
		\draw[->]   (6,2.5) arc(110:-110:5mm) node[midway, left, font=\footnotesize] {$I_2$};
		\draw[->]   (10,2.5) arc(110:-110:5mm) node[midway, left, font=\footnotesize] {$I_3$};
		\draw[->]   (6,6.5)arc(110:-110:5mm) node[midway, left, font=\footnotesize] {$I_4$};
		\draw (n1) node[above left] {$V_2$};
		\draw (n2) node[above right] {$V_3$};
	\end{tikzpicture}
	\caption{}	
	\label{fig:figura2}
\end{figure}
\end{document}


